\documentclass[11pt]{article}

\usepackage{amsmath,amssymb,amsthm,mathtools}
\usepackage[margin=1in]{geometry}

\newtheorem{theorem}{Theorem}[section]
\newtheorem{proposition}[theorem]{Proposition}
\newtheorem{lemma}[theorem]{Lemma}
\newtheorem{remark}[theorem]{Remark}

\newcommand{\kfield}{\Bbbk}
\newcommand{\Uc}{\mathcal U}
\newcommand{\Sc}{\mathcal S}
\newcommand{\fb}{\mathfrak b}
\newcommand{\fB}{\mathfrak B}
\newcommand{\slfr}{\mathfrak{sl}}
\newcommand{\End}{\operatorname{End}}

\begin{document}

\begin{center}
{\Large The punctured bigon and the Drinfeld--Tanisaki double}
\end{center}

\bigskip

\noindent
Throughout let \(\kfield=\mathbb Q(q^{1/n})\), with \(q\) an indeterminate.  We use the simply-connected weight-lattice form of \(\Uc_q(\slfr_n)\), denoted in Klimyk--Schm\"udgen by \(U_q^{\mathrm{ext}}(\slfr_n)\).  Thus one has group-like elements \(\widehat K_1,\ldots,\widehat K_n\), with \(\widehat K_1\cdots \widehat K_n=1\), and Chevalley generators \(E_i,F_i\), \(1\leq i<n\), satisfying
\[
\widehat K_a E_i\widehat K_a^{-1}
=q^{\delta_{a,i}-\delta_{a,i+1}}E_i,\qquad
\widehat K_a F_i\widehat K_a^{-1}
=q^{-\delta_{a,i}+\delta_{a,i+1}}F_i .
\]
Let \(\Uc_q(\fb_+)\) and \(\Uc_q(\fb_-)\) be the Hopf subalgebras generated respectively by the \(\widehat K_a^{\pm1},E_i\) and by the \(\widehat K_a^{\pm1},F_i\).  We use the Drinfeld--Tanisaki skew pairing
\[
\langle \widehat K_i,\widehat K_j\rangle
=q^{1/n-\delta_{ij}},\qquad
\langle E_i,F_j\rangle=\frac{\delta_{ij}}{q^{-1}-q},\qquad
\langle \widehat K_i,F_j\rangle
=\langle E_j,\widehat K_i\rangle=0 .
\]

\section{The \(L\)-operators}

For \(i<j\), let \(E_{ij}\in \Uc_q(\fb_+)\) be the standard positive root vector, normalized by
\[
E_{i,i+1}=E_i,\qquad
E_{ij}=E_{i,j-1}E_{j-1}-q^{-1}E_{j-1}E_{i,j-1}\quad (j>i+1).
\]
For \(i>j\), let \(F_{ij}\in \Uc_q(\fb_-)\) be the corresponding negative root vector, normalized by
\[
F_{j+1,j}=F_j,\qquad
F_{ij}=F_{i-1}F_{i-1,j}-q\,F_{i-1,j}F_{i-1}\quad (i>j+1).
\]
Equivalently, these are Lusztig root vectors for the convex order
\[
\alpha_1,\ \alpha_1+\alpha_2,\ \ldots,\ \alpha_1+\cdots+\alpha_{n-1},\
\alpha_2,\ \alpha_2+\alpha_3,\ \ldots .
\]

Define upper and lower triangular matrices \(L^+=(\ell^+_{ij})\) and \(L^-=(\ell^-_{ij})\) by
\[
\ell^+_{ij}=
\begin{cases}
0,& i>j,\\[2mm]
\widehat K_i,& i=j,\\[2mm]
(q-q^{-1})\,\widehat K_i E_{ij},& i<j,
\end{cases}
\qquad
\ell^-_{ij}=
\begin{cases}
0,& i<j,\\[2mm]
\widehat K_i^{-1},& i=j,\\[2mm]
(q^{-1}-q)\,F_{ij}\widehat K_j^{-1},& i>j .
\end{cases}
\]
These are the FRT \(L\)-operators for the vector representation of \(\Uc_q(\slfr_n)\).

\section{Two standard presentations}

Let \(V=\kfield^n\), with matrix units \(e_{ij}\).  Let \(R\in \End(V\otimes V)\) be the standard type \(A_{n-1}\) Hecke \(R\)-matrix in the normalization compatible with the above pairing:
\[
R
=
q^{-1/n}
\left(
q\sum_i e_{ii}\otimes e_{ii}
+\sum_{i\neq j}e_{ii}\otimes e_{jj}
+(q-q^{-1})\sum_{i<j}e_{ij}\otimes e_{ji}
\right).
\]
For any matrix \(X=(x_{ij})\) with entries in an algebra, write
\[
X_1=\sum_{i,j} e_{ij}\otimes 1\otimes x_{ij},
\qquad
X_2=\sum_{i,j} 1\otimes e_{ij}\otimes x_{ij}.
\]

Let \(\mathcal A_R\) be the algebra generated by entries of an upper triangular matrix \(B=(B_{ij})\) and a lower triangular matrix \(C=(C_{ij})\), subject to
\[
B_{ij}=0\ (i>j),\qquad C_{ij}=0\ (i<j),
\]
the determinant-one relations
\[
B_{11}\cdots B_{nn}=1,\qquad C_{11}\cdots C_{nn}=1,
\]
and the RTT relations
\[
R B_1B_2=B_2B_1R,\qquad
R C_1C_2=C_2C_1R,\qquad
R B_1C_2=C_2B_1R .
\tag{2.1}
\]

\begin{theorem}[RTT presentation of the punctured bigon]
Let \(\fB^*\) be the punctured bigon.  Denote by \(b_{ij}\) the right-going horizontal arc above the puncture, with left state \(i\) and right state \(j\), and by \(c_{ij}\) the right-going horizontal arc below the puncture.  With the standard boundary-state convention, \(b_{ij}=0\) for \(i>j\) and \(c_{ij}=0\) for \(i<j\).  The assignment
\[
B_{ij}\longmapsto b_{ij},\qquad C_{ij}\longmapsto c_{ij}
\]
extends to an algebra isomorphism
\[
\mathcal A_R \xrightarrow{\ \cong\ } \Sc_n(\fB^*).
\]
\end{theorem}

\begin{proof}[Proof/citation]
This is the punctured-bigon specialization of the stated \(\mathrm{SL}_n\)-skein presentation theorem; see L\^e--Sikora, \emph{Stated \(\mathrm{SL}_n\)-skein modules and algebras}, especially the stated basis theorem and the boundary-exchange relations.  The local skein crossing relation gives exactly the RTT equation with the above \(R\)-matrix.  The triangularity \(b_{ij}=0\) for \(i>j\) and \(c_{ij}=0\) for \(i<j\) is the boundary-parallel arc relation.  Finally, the determinant-one relations are the \(\mathrm{SL}_n\) web relation identifying the \(n\)-fold antisymmetrizer with the empty web.  The stated basis theorem gives completeness of these relations.
\end{proof}

\begin{theorem}[FRT presentation of the generalized double]
Let
\[
D=\Uc_q(\fb_-)\bowtie \Uc_q(\fb_+)
\]
be the generalized quantum double with respect to the above Drinfeld--Tanisaki skew pairing.  Write \(x^-=x\bowtie 1\) for \(x\in \Uc_q(\fb_-)\) and \(y^+=1\bowtie y\) for \(y\in \Uc_q(\fb_+)\).  Then the assignment
\[
B_{ij}\longmapsto (\ell^+_{ij})^+,\qquad
C_{ij}\longmapsto (\ell^-_{ij})^-
\]
extends to an algebra isomorphism
\[
\mathcal A_R \xrightarrow{\ \cong\ } D .
\]
\end{theorem}

\begin{proof}[Proof/citation]
This is the standard Faddeev--Reshetikhin--Takhtajan presentation of the Borel halves together with the Drinfeld--Tanisaki cross relation; see Faddeev--Reshetikhin--Takhtajan, Klimyk--Schm\"udgen, and Tanisaki.  The same-sign RTT relations present the two Borel algebras via the \(L^+\)- and \(L^-\)-operators.  The mixed equation
\[
R L^+_1L^-_2=L^-_2L^+_1R
\]
is precisely the defining cross relation in the generalized quantum double associated with the given skew pairing.  The PBW theorem for the double shows that no further relations occur.
\end{proof}

\section{The isomorphism \(\Sc_n(\fB^*)\cong \Uc_q(\fb_-)\bowtie\Uc_q(\fb_+)\)}

Combining the two presentations gives the desired explicit isomorphism.

\begin{theorem}
There is an algebra isomorphism
\[
\Phi:\Sc_n(\fB^*)\xrightarrow{\ \cong\ }
\Uc_q(\fb_-)\bowtie \Uc_q(\fb_+)
\]
given on the horizontal generators by
\[
\Phi(b_{ij})=(\ell^+_{ij})^+,\qquad
\Phi(c_{ij})=(\ell^-_{ij})^- .
\]
Equivalently,
\[
\Phi(b_{ij})=
\begin{cases}
0,& i>j,\\[2mm]
\widehat K_i^+,& i=j,\\[2mm]
(q-q^{-1})\,\widehat K_i^+ E_{ij}^+,& i<j,
\end{cases}
\]
and
\[
\Phi(c_{ij})=
\begin{cases}
0,& i<j,\\[2mm]
(\widehat K_i^-)^{-1},& i=j,\\[2mm]
(q^{-1}-q)\,F_{ij}^-\,(\widehat K_j^-)^{-1},& i>j.
\end{cases}
\]
\end{theorem}

\begin{proof}
Let
\[
\theta_{\mathrm{sk}}:\mathcal A_R\to \Sc_n(\fB^*),
\qquad
\theta_D:\mathcal A_R\to D
\]
be the two isomorphisms from the preceding section.  Then
\[
\Phi=\theta_D\circ \theta_{\mathrm{sk}}^{-1}
\]
is an algebra isomorphism.  The displayed formulas are exactly the images of the generators \(B_{ij}\) and \(C_{ij}\) under \(\theta_D\).
\end{proof}

\section{The quotient giving \(\Uc_q(\slfr_n)\)}

Let \(J\subset D\) be the two-sided ideal generated by
\[
\widehat K_i^+(\widehat K_i^-)^{-1}-1,\qquad 1\leq i\leq n .
\]

\begin{proposition}
There is a canonical algebra isomorphism
\[
D/J \cong \Uc_q(\slfr_n)
\]
sending
\[
E_i^+\mapsto E_i,\qquad
F_i^-\mapsto F_i,\qquad
\widehat K_i^+\mapsto \widehat K_i,\qquad
\widehat K_i^-\mapsto \widehat K_i .
\]
\end{proposition}

\begin{proof}
The natural multiplication map from the two Borel subalgebras into \(\Uc_q(\slfr_n)\) respects the Drinfeld--Tanisaki double relations precisely because the skew pairing is the standard one.  Hence it defines a surjective homomorphism \(D\to \Uc_q(\slfr_n)\).  The elements \(\widehat K_i^+(\widehat K_i^-)^{-1}-1\) lie in its kernel.  By the PBW theorem for the generalized double,
\[
D \cong U_q(\mathfrak n_-)\otimes \kfield[P]\otimes \kfield[P]\otimes U_q(\mathfrak n_+),
\]
where \(P\) is the weight lattice.  Modding out by \(J\) identifies the two torus factors, giving
\[
D/J \cong
U_q(\mathfrak n_-)\otimes \kfield[P]\otimes U_q(\mathfrak n_+),
\]
which is exactly the usual PBW vector-space decomposition of \(\Uc_q(\slfr_n)\).  Therefore the above surjection has kernel exactly \(J\).
\end{proof}

Now let \(I\subset \Sc_n(\fB^*)\) be the two-sided ideal generated by
\[
b_{ii}c_{ii}-1,\qquad 1\leq i\leq n .
\]

\begin{theorem}
There is an algebra isomorphism
\[
\Sc_n(\fB^*)/I \cong \Uc_q(\slfr_n).
\]
Under this isomorphism,
\[
b_{ij}\longmapsto
\begin{cases}
0,& i>j,\\[2mm]
\widehat K_i,& i=j,\\[2mm]
(q-q^{-1})\,\widehat K_iE_{ij},& i<j,
\end{cases}
\]
and
\[
c_{ij}\longmapsto
\begin{cases}
0,& i<j,\\[2mm]
\widehat K_i^{-1},& i=j,\\[2mm]
(q^{-1}-q)\,F_{ij}\widehat K_j^{-1},& i>j.
\end{cases}
\]
Equivalently, the inverse map sends the Chevalley generators to
\[
\widehat K_i\longmapsto [b_{ii}],
\]
\[
E_i\longmapsto (q-q^{-1})^{-1}\,[b_{ii}^{-1}b_{i,i+1}],
\]
\[
F_i\longmapsto (q^{-1}-q)^{-1}\,[c_{i+1,i}c_{ii}^{-1}].
\]
\end{theorem}

\begin{proof}
For each \(i\),
\[
\Phi(b_{ii}c_{ii})
=
\widehat K_i^+(\widehat K_i^-)^{-1}.
\]
Hence \(\Phi(I)=J\).  Therefore \(\Phi\) induces an isomorphism
\[
\Sc_n(\fB^*)/I
\cong
D/J.
\]
By the preceding proposition, \(D/J\cong \Uc_q(\slfr_n)\).  The displayed formulas are obtained by substituting \(\widehat K_i^+=\widehat K_i^-=\widehat K_i\) in the formulas for \(\Phi\).
\end{proof}

\begin{remark}
If one uses the convention in which \(U_q^{\mathrm{ext}}(\slfr_n)\) keeps the central product \(\widehat K_1\cdots \widehat K_n\) instead of imposing \(\widehat K_1\cdots \widehat K_n=1\), then one simply omits the determinant-one relations from the RTT presentation on both sides.  The same proof gives the corresponding extended statement.
\end{remark}

\begin{thebibliography}{9}

\bibitem{FRT}
L.~D. Faddeev, N.~Yu. Reshetikhin and L.~A. Takhtajan,
\emph{Quantization of Lie groups and Lie algebras},
Leningrad Math. J. \textbf{1} (1990), 193--225.

\bibitem{KS}
A.~Klimyk and K.~Schm\"udgen,
\emph{Quantum Groups and Their Representations},
Springer, 1997.

\bibitem{Tanisaki}
T.~Tanisaki,
\emph{Killing forms, Harish-Chandra isomorphisms, and universal \(R\)-matrices for quantum algebras},
in \emph{Infinite Analysis}, Adv. Ser. Math. Phys. \textbf{16},
World Scientific, 1992, 941--961.

\bibitem{CKM}
S.~Cautis, J.~Kamnitzer and S.~Morrison,
\emph{Webs and quantum skew Howe duality},
Math. Ann. \textbf{360} (2014), 351--390.

\bibitem{LeSikora}
T.~T.~Q. L\^e and A.~Sikora,
\emph{Stated \(\mathrm{SL}_n\)-skein modules and algebras},
preprint.

\end{thebibliography}

\end{document}